\documentclass[11pt]{article}

% Match main article formatting
\usepackage[margin=1in]{geometry}
\usepackage{times}

\usepackage[numbers,sort&compress]{natbib}
\usepackage{amsmath,amssymb,amsfonts,amsthm}
\usepackage{graphicx}
\usepackage{url}
\usepackage[english]{babel}

\usepackage{float}
\usepackage{multirow}
\usepackage{listings}
\usepackage{xcolor}
\usepackage{booktabs}
\usepackage[colorlinks=true,linkcolor=blue,citecolor=blue,urlcolor=blue]{hyperref}

\definecolor{codegray}{rgb}{0.5,0.5,0.5}
\definecolor{codeblue}{rgb}{0.0,0.0,0.6}
\definecolor{codegreen}{rgb}{0,0.5,0}

\lstdefinestyle{code}{
  language=Java,
  basicstyle=\fontsize{9}{11}\selectfont\ttfamily,
  numbers=none,
  numberstyle=\tiny\color{codegray},
  stepnumber=1,
  numbersep=5pt,
  backgroundcolor=\color{white},
  showspaces=false,
  showstringspaces=false,
  showtabs=false,
  frame=single,
  rulecolor=\color{black},
  tabsize=2,
  captionpos=b,
  breaklines=true,
  breakatwhitespace=true,
  title=\lstname,
  keywordstyle=\color{codeblue},
  commentstyle=\color{codegreen}\itshape,
  stringstyle=\color{red},
  escapeinside={(*@}{@*)},
  morekeywords={*,...}
}
\lstset{style=code}

\usepackage{textcomp}

\begin{document}

\begin{center}
{\LARGE\bfseries Supplementary Material}
\vspace{0.5em}

{\large Automated Specification Inference via Heuristic Static Analysis\\for Program Comprehension and Test Generation}
\vspace{0.5em}

{\normalsize Brendan Edmonds and Mark Utting}

{\small University of Queensland}
\end{center}

\vspace{1em}
\hrule
\vspace{1em}

\noindent This document provides supplementary material for the main article, including detailed annotation type descriptions, loop invariant heuristic data, a worked inference example, LLM prompt templates, and an approach comparison table.

\tableofcontents
\vspace{1em}

%% ============================================================
\section{Appendix A: Generated Annotation Types}
\label{app:annotation-types}

Unlike traditional JML tools that embed specifications in Javadoc comments, JML-Inferrer generates Java annotations that are compiled into bytecode and queryable at runtime via reflection. Table~\ref{tab:annotation-types} lists the annotation types, all defined in the \texttt{com.jml.inferrer.annotations} package with runtime retention.

\begin{table}[h]
\centering
\caption{Annotation types generated by JML-Inferrer}
\label{tab:annotation-types}
\small
\begin{tabular}{lll}
\toprule
\textbf{Annotation} & \textbf{Target} & \textbf{Description} \\
\midrule
\multicolumn{3}{l}{\textit{Contract Specifications}} \\
\texttt{@Requires} & Method & Precondition (repeatable) \\
\texttt{@Ensures} & Method & Postcondition (repeatable) \\
\texttt{@Signals} & Method & Exceptional postcondition \\
\texttt{@LoopInvariant} & Method & Loop invariant (repeatable) \\
\texttt{@Invariant} & Class & Class invariant (repeatable) \\
\texttt{@Assignable} & Method & Frame condition \\
\midrule
\multicolumn{3}{l}{\textit{Method Classification}} \\
\texttt{@Pure} & Method & No side effects, deterministic \\
\texttt{@Observer} & Method & Reads but doesn't modify state \\
\texttt{@Mutator} & Method & Modifies object state \\
\texttt{@ThreadSafe} & Method/Class & Thread-safe implementation \\
\midrule
\multicolumn{3}{l}{\textit{Nullability}} \\
\texttt{@Nullable} & Field/Param & May be null \\
\texttt{@NonNull} & Field/Param & Never null \\
\midrule
\multicolumn{3}{l}{\textit{Additional Properties}} \\
\texttt{@Immutable} & Field/Class & Immutable value \\
\texttt{@Complexity} & Method & Time/space complexity \\
\texttt{@MustCall} & Method & Required cleanup calls \\
\bottomrule
\end{tabular}
\end{table}

Specifications use JML-like syntax~\cite{leavens2006design, chalin2010jml} with the following constructs:
\begin{itemize}
    \item \texttt{\textbackslash result} -- The method's return value
    \item \texttt{\textbackslash old(expr)} -- Value of \texttt{expr} at method entry
    \item \texttt{\textbackslash nothing} -- Empty frame (no modifications)
    \item \texttt{this.field} -- Instance field reference
    \item \texttt{arr[*]} -- All array elements
    \item \texttt{(\textbackslash forall int k; range; prop)} -- Universal quantification
\end{itemize}

%% ============================================================
\section{Appendix B: Loop Invariant Inference Heuristics}
\label{app:loop-heuristics}

Loop invariant inference is a key challenge in automated verification~\cite{cousot2011contract, wei2011inferring}. Our tool employs a combination of heuristics that have proven effective in practice.

\paragraph{Heuristic 1: Bound-Based Invariants.}
For loops with a clear iteration variable and bound:
\begin{lstlisting}[language=Java]
for (int i = 0; i < n; i++) { ... }
\end{lstlisting}
We infer: $0 \leq i \land i \leq n$.

\paragraph{Heuristic 2: Accumulator Patterns.}
For loops that accumulate values:
\begin{lstlisting}[language=Java]
int sum = 0;
for (int i = 0; i < n; i++) {
    sum += arr[i];
}
\end{lstlisting}
We track the accumulator variable and infer bounds based on symbolic unrolling.

\paragraph{Heuristic 3: Monotonicity Detection.}
We detect whether loop variables are monotonically increasing or decreasing by analyzing the update expression:
\begin{itemize}
    \item $i$\texttt{++} or $i$ \texttt{+= k} (for positive $k$): monotonically increasing
    \item $i$\texttt{--} or $i$ \texttt{-= k} (for positive $k$): monotonically decreasing
\end{itemize}

\paragraph{Heuristic 4: Symbolic Unrolling.}
For loops where pattern-based heuristics fail, we perform bounded symbolic unrolling (default: 3 iterations) to detect emergent patterns. If a consistent relationship is observed, it is generalized to an invariant candidate.

\paragraph{Heuristic 5: While-Loop Counter Detection.}
For \texttt{while} loops, which lack explicit initialization and update expressions, the tool scans the loop body for counter variables---identified by increment/decrement operations (\texttt{i++}, \texttt{i-{}-}, \texttt{i += n}) and compound assignments. Detected counters are then combined with bounds extracted from the loop condition (e.g., \texttt{while (i < n)} yields $i \geq 0$ and $i < n$). The same accumulator and variable-relationship analyses used for \texttt{for} loops are applied to the body. Common iterator idioms such as \texttt{while (iter.hasNext())} are also recognized.

\paragraph{Heuristic 6: Enhanced For-Each Analysis.}
For enhanced \texttt{for-each} loops, the tool infers element membership invariants (\texttt{elem} is an element of the iterable) and collection-size bounds (processed elements $\leq$ collection size). Variable-relationship analysis (e.g., max/min tracking) is applied to the loop body, producing richer invariants than element membership alone.

\paragraph{Fallback Behavior.}
When all heuristics fail, the tool produces a conservative postcondition that does not constrain the loop's effect. We record this as a ``weak specification'' for analysis.

Table~\ref{tab:loop-heuristics} reports the success rate of each heuristic on the 68 loops encountered across 48 looping methods in our Commons Lang evaluation dataset.

\begin{table}[h]
\centering
\caption{Loop invariant inference success rates}
\label{tab:loop-heuristics}
\begin{tabular}{lrr}
\toprule
\textbf{Heuristic} & \textbf{Loops} & \textbf{Success Rate} \\
\midrule
Bound-based (for) & 32 & 93.8\% \\
Accumulator pattern & 15 & 80.0\% \\
Monotonicity & 13 & 84.6\% \\
Symbolic unrolling & 8 & 62.5\% \\
While-loop counter & --- & --- \\
Enhanced for-each & --- & --- \\
\midrule
Overall & 68 & 85.3\% \\
Fallback (weak spec) & 10 & --- \\
\bottomrule
\end{tabular}
\end{table}

%% ============================================================
\section{Appendix C: Worked Example --- End-to-End Inference}
\label{app:worked-example}

This appendix presents a complete example showing how JML-Inferrer analyzes a method and derives its specification.

\subsection{Input Method}

Consider the following utility method that validates and returns an array element:

\begin{lstlisting}[language=Java, caption={Input method before specification inference}, label={lst:before-supp}]
public static String getElement(String[] arr, int index) {
    if (arr == null) {
        throw new IllegalArgumentException("Array cannot be null");
    }
    if (index < 0 || index >= arr.length) {
        throw new IndexOutOfBoundsException("Invalid index");
    }
    return arr[index];
}
\end{lstlisting}

\subsection{Precondition Inference (WP-Inspired)}

The tool identifies guard-and-throw patterns that correspond to preconditions:

\begin{enumerate}
    \item \textbf{First if-statement}: The pattern \texttt{if (arr == null) throw ...} is recognized as a null-check guard, yielding precondition \texttt{arr != null}.

    \item \textbf{Second if-statement}: The pattern \texttt{if (index < 0 || index >= arr.length) throw ...} is recognized as a range-check guard, yielding preconditions \texttt{index >= 0} and \texttt{index < arr.length}.
\end{enumerate}

These heuristic inferences correspond to what a formal WP analysis would derive---the weakest precondition for the return statement \texttt{return arr[index]} with postcondition \texttt{true} is exactly $\texttt{arr} \neq \texttt{null} \land 0 \leq \texttt{index} < \texttt{arr.length}$.

\noindent\textbf{Inferred Preconditions}: \texttt{arr != null}, \texttt{index >= 0}, \texttt{index < arr.length}

\subsection{Postcondition Inference (SP-Inspired)}

The tool analyzes the return statement to derive postconditions:

\begin{enumerate}
    \item The return expression \texttt{arr[index]} is an array access on input parameters, yielding $\texttt{\textbackslash result} = \texttt{arr[index]}$.
    \item Since no fields are modified, the method is classified as \texttt{@Pure}.
\end{enumerate}

\noindent\textbf{Inferred Postcondition}: \texttt{\textbackslash result == arr[index]}

\subsection{Output: Annotated Method}

\begin{lstlisting}[language=Java, caption={Output method with inferred JML specifications}, label={lst:after-supp}]
@Requires("arr != null")
@Requires("index >= 0")
@Requires("index < arr.length")
@Ensures("\\result == arr[index]")
@Assignable("\\nothing")
@Pure
@Complexity(time = "O(1)", space = "O(1)")
public static String getElement(String[] arr, int index) {
    if (arr == null) {
        throw new IllegalArgumentException("Array cannot be null");
    }
    if (index < 0 || index >= arr.length) {
        throw new IndexOutOfBoundsException("Invalid index");
    }
    return arr[index];
}
\end{lstlisting}

The tool also infers that the method is \texttt{@Pure} (no side effects) and has constant time complexity.

%% ============================================================
\section{Appendix D: LLM Prompt Templates}
\label{app:prompts}

The following prompts were used for LLM-based test generation in each experimental phase. Prompts were kept minimal, differing only in the inclusion of specifications or source code.

\subsection{Phase 1 Prompt (Signature Only)}
\begin{quote}
\textit{Write unit tests for the following function signature under typical development constraints. Provide the tests as a developer would normally write them, without any formal specification guidance:}

\texttt{[method signature]}
\end{quote}

\subsection{Phase 2 Prompt (Signature + Guidance)}
\begin{quote}
\textit{Given the following function signature, write a comprehensive set of unit tests. Ensure that the tests cover normal behavior, edge cases, and potential failure scenarios:}

\texttt{[method signature]}
\end{quote}

\subsection{Phase 3 Prompt (Signature + Specification)}
\begin{quote}
\textit{Given the following function signature and formal specification, write a comprehensive set of unit tests. Ensure that the tests cover normal behavior, edge cases, and potential failure scenarios. The specification uses JML notation where @requires specifies preconditions and @ensures specifies postconditions:}

\texttt{[method signature with JML specification]}
\end{quote}

\subsection{Phase 4 Prompt (Source Code)}
\begin{quote}
\textit{Given the following function implementation, write a comprehensive set of unit tests. Ensure that the tests cover normal behavior, edge cases, and potential failure scenarios:}

\texttt{[full method source code]}
\end{quote}

%% ============================================================
\section{Appendix E: Approach Comparison}
\label{app:comparison}

Table~\ref{tab:approach-comparison-supp} summarizes the key differences between our approach and related specification inference methods.

\begin{table}[h]
\centering
\caption{Comparison of specification inference approaches}
\label{tab:approach-comparison-supp}
\small
\begin{tabular}{lcccccc}
\toprule
\textbf{Approach} & \textbf{Input} & \textbf{Technique} & \textbf{Lang.} & \textbf{Prec.} & \textbf{Rec.} & \textbf{Cov.} \\
\midrule
Our tool & Source & Heuristic static & Java & 94.2\% & 89.3\% & 99.0\% \\
Jdoctor & Javadoc & NLP & Java & 92\% & 83\% & 45.2\% \\
PreInfer & Execution & Symbolic & C\# & N/R & N/R & N/R \\
Daikon & Traces & Dynamic & Multi & varies & varies & varies \\
LLM-based & Source & Neural & Multi & 81.3\% & 88.7\% & 100\% \\
\bottomrule
\end{tabular}
\end{table}

Key observations:
\begin{itemize}
    \item Our approach achieves the highest precision among automated methods, due to its conservative heuristic pattern matching grounded in common code idioms.
    \item LLM-based inference achieves higher recall but lower precision, and lacks determinism.
    \item Jdoctor achieves high precision but is limited by documentation availability.
    \item Dynamic approaches like Daikon require execution traces and may miss corner cases not exercised during testing.
\end{itemize}

\bibliographystyle{plainnat}
\bibliography{references}

\end{document}
